\documentclass[12pt]{article}
\usepackage[margin=0.6in,top=1.15in,bottom=0.55in]{geometry}
\usepackage{lmodern}
\usepackage{microtype}
\usepackage{titlesec}
\usepackage{enumitem}
\usepackage[hidelinks]{hyperref}
\usepackage{../../latex/grader-worksheet}
\setlength{\parindent}{0pt}
\setlength{\parskip}{0.25em}
\titleformat{\section}{\large\bfseries\scshape}{}{0pt}{}
\GraderSetup{assignment-id=U1L19LL,template-revision=1}
\begin{document}
\begin{center}
{\LARGE\bfseries Final Portfolio Defense}\par\vspace{0.05em}
{\large\scshape What It Means to Think Straight}\par\vspace{0.1em}
\rule{0.78\textwidth}{0.5pt}\par\vspace{0.15em}
\end{center}
\noindent\textbf{Purpose:} Construct and test a defense of intellectual growth using Whately, at least three portfolio artifacts, one revision, one Shakespeare application, and one concrete transfer.
\vspace{0.3em}

\section*{Part 1: Definition and Thesis}
\textbf{Question 1.1}\\[-0.15em]
Define \emph{disciplined thinker} through observable habits. Include at least three accurate Whately habits and one limit on what logic can do.
\GraderAnswerBox{Part 1 Q1}{2.55in}

\textbf{Question 1.2}\\[-0.15em]
Write the controlling claim of your defense. State a specific pattern of growth that your portfolio can prove; do not claim perfection or mere completion.
\GraderAnswerBox{Part 1 Q2}{2.45in}

\newpage
\section*{Part 2: Artifact Evidence}
\textbf{Question 2.1}\\[-0.15em]
Present Artifact 1. Identify it, describe or quote the relevant reasoning, name the logical habit, and explain how it supports your thesis.
\GraderAnswerBox{Part 2 Q1}{2.7in}

\textbf{Question 2.2}\\[-0.15em]
Present Artifact 2 using the same evidence-and-commentary standard. This artifact must show a different stage or kind of reasoning from Artifact 1.
\GraderAnswerBox{Part 2 Q2}{2.7in}

\newpage
\section*{Part 2 Continued: Required Artifacts}
\textbf{Question 2.3}\\[-0.15em]
Present Artifact 3. Across Questions 2.1--2.3, at least one artifact must be substantively revised and at least one must use \emph{Macbeth} or \emph{Julius Caesar}. Identify those requirements explicitly.
\GraderAnswerBox{Part 2 Q3}{3.0in}

\textbf{Question 2.4}\\[-0.15em]
State the warrant connecting each artifact to the pattern in your thesis. Explain why these are evidence of growth rather than isolated correct answers.
\GraderAnswerBox{Part 2 Q4}{2.55in}

\newpage
\section*{Part 3: Test the Defense}
\textbf{Question 3.1}\\[-0.15em]
Name one counterpressure: a remaining weakness, contrary artifact, or fair objection to your thesis. Then qualify the thesis so it remains accurate without becoming empty.
\GraderAnswerBox{Part 3 Q1}{2.7in}

\textbf{Question 3.2}\\[-0.15em]
Use the Shakespeare artifact as a judgment mirror. Identify the character's pressure point and explain which disciplined habit the passage makes necessary; avoid plot summary and moral boasting.
\GraderAnswerBox{Part 3 Q2}{2.7in}

\newpage
\section*{Part 4: Transfer and Closing Defense}
\textbf{Question 4.1}\\[-0.15em]
Make one concrete transfer claim. Name a setting beyond this unit, a likely reasoning pressure, and the exact action you will take because of one Whately habit.
\GraderAnswerBox{Part 4 Q1}{2.45in}

\textbf{Question 4.2}\\[-0.15em]
Write the closing paragraph or oral-defense notes. Restate the qualified thesis, synthesize the three artifacts, and answer: What does it mean to think straight?
\GraderAnswerBox{Part 4 Q2}{3.05in}

\newpage
\section*{Part 4 Continued: Final Logic Check}
\textbf{Question 4.3}\\[-0.15em]
Audit your own defense. Identify one place where terms, evidence, inference, source judgment, rhetoric, or qualification still needs attention, and write the exact revision you will make.
\GraderAnswerBox{Part 4 Q3}{4.85in}

\newpage
\section*{Part 5: Review}
\textbf{Question 5.1}\\[-0.15em]
An early portfolio caption says, ``I changed my answer, so this proves I became a disciplined thinker.'' Use the Week 18 revision standard to name what evidence the caption still needs.
\GraderAnswerBox{Part 5 Q1}{2.35in}

\textbf{Question 5.2}\\[-0.15em]
A debater admits privately, ``I chose the policy first,'' but tells the audience, ``Only an enemy of safety would question this plan.'' Use a prior-week tool to identify the pressure point and one narrower honest claim.
\GraderAnswerBox{Part 5 Q2}{2.45in}
\end{document}
